% SPDX-License-Identifier: LPPL-1.3c
% Copyright (C) 2026 Christophe Jorssen
% Author and Current Maintainer: Christophe Jorssen <christophe.jorssen@gmail.com>
% LPPL maintenance status: maintained
% This file is part of luafemm. See LICENSE and LICENSES.md for its terms.
% Build from the repository root with make manual. Examples execute on LuaTeX.
% The PGF documentation environments also used by tikz-ext provide indexed
% keys, library/command signatures and adjacent source/rendered examples.
\documentclass[a4paper,11pt]{ltxdoc}
\usepackage[margin=23mm]{geometry}
\usepackage{fontspec}
\usepackage{amsmath,amssymb,booktabs,longtable,multicol}
\usepackage{xkeyval,calc,listings,fp,pifont,xxcolor}
\usepackage{tikz,pgfplots}
\usepackage{makeidx}
\makeindex
\usepackage[hyperindex=false]{hyperref}
\usepackage{xurl}
\hypersetup{colorlinks=true,linkcolor=blue!55!black,urlcolor=blue!55!black,
  pdftitle={The luafemm TikZ and PGFPlots libraries},pdfauthor={Christophe Jorssen}}
\input{pgfmanual-en-macros.tex}
\input{doc/manual-examples.tex}
\usetikzlibrary{femm,calc}
\usepgfplotslibrary{femm}
\pgfplotsset{compat=1.18}
\tracinglostchars=3
\setlength\emergencystretch{2em}
\title{The \textsf{luafemm} TikZ and PGFPlots libraries\\
  \large Generic LuaTeX magnetostatics and field profiles}
\author{Christophe Jorssen\\
  \normalsize\href{mailto:christophe.jorssen@gmail.com}{\nolinkurl{christophe.jorssen@gmail.com}}}
\date{Version \luafemmversion\quad---\quad\luafemmversiondate}
\begin{document}
\maketitle
\begin{abstract}
Declare magnetic media on ordinary TikZ paths, solve a planar finite-element
problem entirely in Lua, and draw oriented field lines or PGFPlots profiles.
The same generic implementation works with plain LuaTeX, LuaLaTeX and
ConTeXt MkIV. This manual documents every supported public key, the numerical
conventions, and the limitations of this development version.
\end{abstract}
\input{doc/sections/project.tex}
\clearpage
\tableofcontents
\clearpage
\input{doc/sections/getting-started.tex}
\input{doc/sections/tutorial.tex}
\input{doc/sections/model.tex}
\input{doc/sections/mesh-quality.tex}
\input{doc/sections/domains.tex}
\input{doc/sections/boundaries.tex}
\input{doc/sections/components.tex}
\input{doc/sections/materials.tex}
\input{doc/sections/field.tex}
\input{doc/sections/cache.tex}
\input{doc/sections/profiles.tex}
\input{doc/sections/numerics.tex}
\input{doc/sections/development.tex}
\appendix
\input{doc/sections/legacy.tex}
\input{doc/sections/format-examples.tex}
\section{Complete FEMM material catalogue}
\label{sec:catalogue}
The identifiers below are exact choices for |material| and |library|.
The original names are retained, including the two distinct Supermalloy
records. The full data and original folder paths are available through
the Lua catalogue and in \texttt{docs/materials-catalog.md}.
\input{build/doc/materials-catalog.tex}
\clearpage
\printindex
\end{document}
